\documentclass[12pt]{article}
\usepackage[utf8]{inputenc}
\usepackage[T1]{fontenc}
\usepackage{hyperref}
\title{Saaroa Language Resources for K-12 Curriculum Development}
\author{Compiled by Claude Sonnet 5.5 (claude-sonnet-5-5) for Andr\'as Kornai}
\date{}
\begin{document}
\maketitle

\section{Introduction}
Saaroa (Lha'alua; ISO 639-3 sxr) is a critically endangered Austronesian language of Taiwan, with fewer than about 10 speakers reported. The community numbers about 400 and lives in Taoyuan and Kaochung villages (Taoyuan District, Kaohsiung), with Maya village in Namasia District mentioned by one source. Documentation exists, but no school curriculum was found.

\section{Primary Educational Resources}
No Saaroa primers or textbooks were identified in the searches. One blog says many learning resources are in Traditional Chinese and that a central site covers words, dialogues and audio, but the site was not named, so it is not cited here.

\section{Linguistic and Reference Materials}
\begin{itemize}
\item Wikipedia, ``Saaroa language'': \url{https://en.wikipedia.org/wiki/Saaroa_language}
\item Omniglot, Saaroa writing and orthography (states a standard orthography was set by Taiwan's Council of Indigenous Peoples and Ministry of Education): \url{https://omniglot.com/writing/saaroa.htm}
\item The Language Closet, ``Languages of Taiwan: Saaroa (Lha'alua)'': \url{https://thelanguagecloset.com/2021/12/04/languages-of-taiwan-saaroa-lhaalua/}
\item Numerals of Saaroa, Max Planck EVA: \url{https://lingweb.eva.mpg.de/channumerals/Saaroa.htm}
\end{itemize}

\section{Religious and Community Resources}
None specific to Saaroa were found.

\section{Digital and Audio Resources}
\begin{itemize}
\item ELDP project ``A Documentation of Saaroa, a Moribund Austronesian Language of Taiwan'' (Chia-jung Pan; texts including conversation, stories and ritual speech): \url{https://cultureincrisis.org/projects/a-documentation-of-saaroa-a-moribund-austronesian-language-of-taiwan}
\item Archive record for the same collection: \url{https://eldp.access.preservica.com/dk0379/}
\item Omniglot, counting in Saaroa (with some audio): \url{https://omniglot.com/language/numbers/saaroa.htm}
\end{itemize}

\section{Curriculum Development Context}
With so few speakers, a curriculum would depend on the archived documentation corpus and the official orthography, plus work with the Saaroa community and Taiwanese indigenous-language education authorities. The archive contents were not examined in detail.

\end{document}